\section{Discussion: implications for predictive modelling}
\label{sec:discussion}
\subsection{Conservation, resolution and experimental agreement}
The results separate three requirements that a deposition calculation
must satisfy. Element balance tests the consistency of the discrete
transport--equilibrium--wall coupling. Resolution studies test how much
the computed observables depend on numerical choices. Experimental
comparison tests the combined consequences of that numerical model,
its physical assumptions and its input data. Passing the first test
does not establish the other two.

This distinction is consequential here. The raw material residues are
small, but the original upwind fine/medium profile differences are
2.06--10.58\%.
Meanwhile, total deposition changes by less than 0.014\% under the same
refinement. An integral balance can therefore appear stable while the
spatial distribution remains sensitive. A binned peak introduces an
additional loss of information: nearly equal neighbouring bins can
exchange rank and produce a large change in the reported peak
temperature. Profile differences, centroids and deposited amounts
should accompany peak comparisons when assessing this model.
The production scheme and mesh study reduces the most sensitive
limited-linear profile change to \SteelLLThree\% between 41,472 and
93,312 cells, without a consistent apparent order.
The saved-profile aggregation check in \cref{sec:bin-audit} strengthens
this distinction: wider bins can suppress the apparent mesh difference,
but the result depends strongly on their origin. Conversely, the small
binned transport distances support stability of broad deposition location
between the tested meshes without qualifying local peak heights.

The physical-input variations also need to be interpreted against this
resolution evidence. Their ranges are finite responses over selected
cases, not components of a probabilistic uncertainty budget. They cannot
be added in quadrature or used to infer a universal parameter ranking.
The refined BiI perturbation retains a large change in iodide partition,
whereas several transport and accommodation responses are comparable
in magnitude to the unresolved spatial differences. The latter require
further numerical qualification before small physical effects can be
distinguished reliably.

\subsection{What the experimental comparisons can identify}
The source constructions expose a limitation of inferring release
chemistry from final deposits. Q1 uses deposit-derived compound masses
for LBE, so agreement in the iodide split partly reflects information
already supplied to the calculation. Q2 allows the split to emerge,
but combines a deposit-derived iodine loading with an unmeasured metal
saturation. The response to BiI thermodynamics then shows that the
same observed split may constrain several coupled assumptions rather
than identify a unique source or thermodynamic description.

The silica results reveal a separate ambiguity. Both the printed peak
temperatures and the digitised spatial scan appear plausible when
considered individually, yet they disagree through the reconstructed
temperature-coordinate mapping. The two diagnostic shifts substantially
improve the profiles without resolving which mapping is physically
correct; their overlap ranking reverses between upwind and limited-linear
advection. Moreover, these shifts use labels from the same experiment
being compared. Although they were not fitted to CFD output, their
improved agreement is still not an independent validation test.

The untuned single-species comparison with Lobresco et al.\
\cite{Lobresco2026assessment} supplies an additional experimental test:
the onset-temperature RMS deviation is \LobRMS{} K, but the computed
peaks lie \LobShiftMin--\LobShiftMax{} cm upstream and the shutdown
transient and non-dilute source region remain unrepresented. It strengthens
the evidence for the transport and wall law within those limits.

The practical contribution is therefore a framework for testing explicit
input hypotheses while retaining conservation and exposing profile
discrepancies. The present data do not uniquely separate release,
thermodynamic, coordinate and omitted-chemistry effects. Fitting an
interfacial coefficient to a single peak under these conditions could
compensate for an error elsewhere in the model.

\subsection{Evidence needed to extend the present scope}
The scheme and mesh study now extends the shifted-BiI steel branch to
93,312 cells with bounded advection, and full-duration corrector checks
are available for three cases. Spatial qualification still needs to extend
beyond that one refined branch and include cases with competing binned
maxima. Full profiles and integral observables should both be assessed.
The full-duration shifted-BiI corrector comparison remains unreported.

On the experimental side, the most direct constraints would be the
original thermocouple and scan coordinates, controller reference
conditions and an independently measured release history. Qualified
BiI thermodynamics and the provenance of the constructed condensed
records are also needed before selecting a thermochemical branch on
physical grounds. These requirements address specific ambiguities
demonstrated by the completed cases.

Finally, comparison with post-experiment deposits requires the release
and cooling protocol to be represented or its effect independently
bounded. The current 60 s calculations preserve the assumed three-hour
mean source rate, while the flow-stop restarts retain the original
temperature field. Neither calculation supplies the missing history.
The role of potassium chemistry, substrate reactions, finite gas
kinetics and boat obstruction remains outside the evidence presented
here and cannot be inferred from the present sensitivity ranges.
